\QHeading{第1题　给图形涂色表示它下面的小数。}

\par\noindent\makebox[\linewidth][c]{
\begin{tikzpicture}[x=1bp,y=-1bp,color=black,draw=black,text=black]
\begin{scope}\clip (56.0,115.0) rectangle (533.0,266.0);
\hdhead{56}{小数的意义与性质——复习}
\hdsection{1}{137}{给图形涂色表示它下面的小数。}
\begin{scope}[scale=.6666667,draw=black,text=black]\draw[hdgrid] (155,182) rectangle (296,329);\foreach \i in {1,...,9}{\draw[hdgrid] ({155+14.1*\i},182)--({155+14.1*\i},329);}\foreach \j in {1,...,0}{\draw[hdgrid] (155,{182+147.0*\j})--(296,{182+147.0*\j});}\end{scope}
\begin{scope}[scale=.6666667,draw=black,text=black]\node[inner sep=0pt,font=\normalsize,text=black] at (133.000,192.000) {①};\end{scope}
\begin{scope}[scale=.6666667,draw=black,text=black]\node[inner sep=0pt,font=\normalsize,text=black] at (225.500,354.000) {0.6};\end{scope}
\begin{scope}[scale=.6666667,draw=black,text=black]\draw[hdgrid] (399,182) rectangle (541,329);\foreach \i in {1,...,9}{\draw[hdgrid] ({399+14.2*\i},182)--({399+14.2*\i},329);}\foreach \j in {1,...,9}{\draw[hdgrid] (399,{182+14.7*\j})--(541,{182+14.7*\j});}\end{scope}
\begin{scope}[scale=.6666667,draw=black,text=black]\node[inner sep=0pt,font=\normalsize,text=black] at (377.000,192.000) {②};\end{scope}
\begin{scope}[scale=.6666667,draw=black,text=black]\node[inner sep=0pt,font=\normalsize,text=black] at (470.000,354.000) {0.97};\end{scope}
\begin{scope}[scale=.6666667,draw=black,text=black]\draw[hdgrid] (645,182) rectangle (791,329);\foreach \i in {1,...,9}{\draw[hdgrid] ({645+14.6*\i},182)--({645+14.6*\i},329);}\foreach \j in {1,...,0}{\draw[hdgrid] (645,{182+147.0*\j})--(791,{182+147.0*\j});}\end{scope}
\begin{scope}[scale=.6666667,draw=black,text=black]\node[inner sep=0pt,font=\normalsize,text=black] at (623.000,192.000) {③};\end{scope}
\begin{scope}[scale=.6666667,draw=black,text=black]\node[inner sep=0pt,font=\normalsize,text=black] at (718.000,354.000) {0.2};\end{scope}
\hdsection{2}{402}{化简下面的小数。}
\q{124}{448}{180}{13}{$0.40=$}
\q{124}{492}{180}{13}{$8.80=$}
\q{124}{536}{180}{13}{$2.030=$}
\q{313}{448}{180}{13}{$1.50=$}
\q{313}{492}{180}{13}{$6.030=$}
\q{313}{536}{180}{13}{$3.0600=$}
\q{501}{448}{180}{13}{$3.200=$}
\q{501}{492}{180}{13}{$12.400=$}
\q{501}{536}{180}{13}{$109.200=$}
\q{693}{448}{180}{13}{$10.050=$}
\q{693}{492}{180}{13}{$0.900=$}
\q{693}{536}{180}{13}{$0.070=$}
\hdsection{3}{591}{排一排。}
\q{124}{640}{30}{13}{①}
\q{154}{640}{91}{13}{2.3}
\q{210}{640}{91}{13}{2.03}
\q{277}{640}{91}{13}{23.2}
\q{345}{640}{91}{13}{23.23}
\q{422}{640}{91}{13}{2.303}
\q{124}{726}{30}{13}{②}
\q{154}{726}{91}{13}{0.11}
\q{210}{726}{91}{13}{0.1}
\q{277}{726}{91}{13}{0.111}
\q{345}{726}{91}{13}{0.101}
\q{422}{726}{91}{13}{1.101}
\q{153}{683}{28}{12.5}{（}
\q{209}{683}{25}{12.5}{）}
\q{244}{683}{28}{12.5}{（}
\q{301}{683}{25}{12.5}{）}
\q{337}{683}{28}{12.5}{（}
\q{392}{683}{25}{12.5}{）}
\q{430}{683}{28}{12.5}{（}
\q{486}{683}{25}{12.5}{）}
\q{525}{683}{28}{12.5}{（}
\q{578}{683}{25}{12.5}{）}
\q{224}{683}{28}{13}{$<$}
\q{316}{683}{28}{13}{$<$}
\q{408}{683}{28}{13}{$<$}
\q{501}{683}{28}{13}{$<$}
\q{153}{769}{28}{12.5}{（}
\q{209}{769}{25}{12.5}{）}
\q{244}{769}{28}{12.5}{（}
\q{301}{769}{25}{12.5}{）}
\q{337}{769}{28}{12.5}{（}
\q{392}{769}{25}{12.5}{）}
\q{430}{769}{28}{12.5}{（}
\q{486}{769}{25}{12.5}{）}
\q{525}{769}{28}{12.5}{（}
\q{578}{769}{25}{12.5}{）}
\q{224}{769}{28}{13}{$>$}
\q{316}{769}{28}{13}{$>$}
\q{408}{769}{28}{13}{$>$}
\q{501}{769}{28}{13}{$>$}
\hdsection{4}{825}{$\square$里最大能填几?}
\q{122}{879}{162}{13}{$0.23>0.$}
\q{249}{879}{162}{13}{$2$}
\q{371}{879}{162}{13}{$0.1$}
\q{444}{879}{162}{13}{$<0.15$}
\q{627}{879}{162}{13}{$3.4$}
\q{704}{879}{162}{13}{$<3.42$}
\q{122}{937}{162}{13}{$2.56>2.5$}
\q{371}{937}{162}{13}{$45.6>4$}
\q{494}{937}{162}{13}{$.6$}
\q{627}{937}{162}{13}{$7.31>7.$}
\q{761}{937}{162}{13}{$1$}
\begin{scope}[scale=.6666667,draw=black,text=black]\draw[line width=.5pt] (207,869) rectangle (242,904);\end{scope}
\begin{scope}[scale=.6666667,draw=black,text=black]\draw[line width=.5pt] (410,869) rectangle (445,904);\end{scope}
\begin{scope}[scale=.6666667,draw=black,text=black]\draw[line width=.5pt] (666,869) rectangle (701,904);\end{scope}
\begin{scope}[scale=.6666667,draw=black,text=black]\draw[line width=.5pt] (224,927) rectangle (259,962);\end{scope}
\begin{scope}[scale=.6666667,draw=black,text=black]\draw[line width=.5pt] (456,927) rectangle (491,962);\end{scope}
\begin{scope}[scale=.6666667,draw=black,text=black]\draw[line width=.5pt] (717,927) rectangle (752,962);\end{scope}
\hdsection{5}{997}{小数改写。}
\q{124}{1045}{350}{12.5}{① 把0.6改写成两位小数。}
\q{501}{1045}{339}{12.5}{② 把2.5改写成三位小数。}

\begin{scope}[color=answerred,draw=answerred,text=answerred]\fill[answerred,opacity=.28](103.333,121.333) rectangle (159.733,219.333);
\fill[answerred,opacity=.28](266.000,121.333) rectangle (351.200,219.333);
\fill[answerred,opacity=.28](351.200,121.333) rectangle (360.667,189.933);
\fill[answerred,opacity=.28](430.000,121.333) rectangle (449.467,219.333);
\ans{124.00}{298.00}{61.33}{9.5}{$\vphantom{0}0.4$}
\ans{124.00}{327.33}{61.33}{9.5}{$\vphantom{0}8.8$}
\ans{132.67}{356.67}{61.33}{9.5}{$\vphantom{0}2.03$}
\ans{250.67}{298.00}{61.33}{9.5}{$\vphantom{0}1.5$}
\ans{258.67}{327.33}{61.33}{9.5}{$\vphantom{0}6.03$}
\ans{266.00}{356.67}{61.33}{9.5}{$\vphantom{0}3.06$}
\ans{385.33}{298.00}{61.33}{9.5}{$\vphantom{0}3.2$}
\ans{394.00}{327.33}{61.33}{9.5}{$\vphantom{0}12.4$}
\ans{401.33}{356.67}{61.33}{9.5}{$\vphantom{0}109.2$}
\ans{522.00}{298.00}{61.33}{9.5}{$\vphantom{0}10.05$}
\ans{514.67}{327.33}{61.33}{9.5}{$\vphantom{0}0.9$}
\ans{514.67}{356.67}{61.33}{9.5}{$\vphantom{0}0.07$}
\ans{116.00}{454.00}{62.67}{10}{$\vphantom{0}2.03$}
\ans{178.00}{454.00}{62.67}{10}{$\vphantom{0}2.3$}
\ans{236.67}{454.00}{62.67}{10}{$\vphantom{0}2.303$}
\ans{300.67}{454.00}{62.67}{10}{$\vphantom{0}23.2$}
\ans{363.33}{454.00}{62.67}{10}{$\vphantom{0}23.23$}
\ans{116.00}{510.00}{62.67}{10}{$\vphantom{0}1.101$}
\ans{176.67}{510.00}{62.67}{10}{$\vphantom{0}0.111$}
\ans{239.33}{510.00}{62.67}{10}{$\vphantom{0}0.11$}
\ans{300.67}{510.00}{62.67}{10}{$\vphantom{0}0.101$}
\ans{363.33}{510.00}{62.67}{10}{$\vphantom{0}0.1$}
\ans{144.00}{584.00}{33.33}{10}{$\vphantom{0}2$}
\ans{279.33}{584.00}{33.33}{10}{$\vphantom{0}4$}
\ans{450.67}{584.00}{33.33}{10}{$\vphantom{0}1$}
\ans{155.33}{622.67}{33.33}{10}{$\vphantom{0}5$}
\ans{310.00}{622.67}{33.33}{10}{$\vphantom{0}4$}
\ans{484.00}{622.67}{33.33}{10}{$\vphantom{0}2$}
\ans{95.33}{733.33}{180.00}{11}{$\vphantom{0}0.6=0.60$}
\ans{348.00}{733.33}{180.00}{11}{$\vphantom{0}2.5=2.500$}
\end{scope}
\end{scope}\pgfresetboundingbox\path[use as bounding box] (56.0,115.0) rectangle (533.0,266.0);
\end{tikzpicture}}\par
\par\noindent




\par





\par\noindent

\QHeading{第1题　涂色不必照抄同一位置}

\Answer{0.6可以涂10份中的任意6份；0.97可以涂100份中的97份，也可以先涂满再留3份空白；0.2可以涂10份中的任意2份。核对的是同一个整体所占的份数，不是涂色位置。}

\QHeading{第2题　化简下面的小数。}

\begin{oralgrid}
\OralItem{1}{0.40}{0.4}
\OralItem{2}{1.50}{1.5}
\OralItem{3}{3.200}{3.2}
\OralItem{4}{10.050}{10.05}
\OralItem{5}{8.80}{8.8}
\OralItem{6}{6.030}{6.03}
\OralItem{7}{12.400}{12.4}
\OralItem{8}{0.900}{0.9}
\OralItem{9}{2.030}{2.03}
\OralItem{10}{3.0600}{3.06}
\OralItem{11}{109.200}{109.2}
\OralItem{12}{0.070}{0.07}\end{oralgrid}
\par

\par\noindent

\QHeading{第3题　排一排。}

\PlainItem{①2.32.0323.223.232.303}
\par

\PlainItem{（{\color{answerred}$2.03$}）$<$（{\color{answerred}$2.3$}）$<$（{\color{answerred}$2.303$}）$<$（{\color{answerred}$23.2$}）$<$（{\color{answerred}$23.23$}）}
\PlainItem{②0.110.10.1110.1011.101}
\PlainItem{（{\color{answerred}$1.101$}）$>$（{\color{answerred}$0.111$}）$>$（{\color{answerred}$0.11$}）$>$（{\color{answerred}$0.101$}）$>$（{\color{answerred}$0.1$}）}
\par\noindent

\QHeading{第4题　$\square$里最大能填几?}

\Equation{0.23>0.\textcolor{answerred}{2}20.1\textcolor{answerred}{4}<0.153.4\textcolor{answerred}{1}<3.42}
\par

\Equation{2.56>2.5\textcolor{answerred}{5}45.6>4\textcolor{answerred}{4}.67.31>7.\textcolor{answerred}{2}1}
\par\noindent

\QHeading{第5题　小数改写。}

\PlainItem{① 把0.6改写成两位小数。② 把2.5改写成三位小数。}
\par

\Equation{\textcolor{answerred}{0.6=0.60}\textcolor{answerred}{2.5=2.500}}

